\honest{Joint Configurations}{Eliezer Yudkowsky}{2008}

The key to quantum mechanics is to realize ``on a truly gut level that configurations are about more than one particle.''

I send two photons at once, one from B and one from C, into the half-silvered mirror D. Deflection multiplies the amplitude by $i$, going straight by 1. Of the four routes, two end in the same configuration, one photon going to each detector, and their amplitudes are $1$ and $-1$. They cancel. So the two detectors never fire one each; both photons go to Detector 1 or both to Detector 2, equally often. ``I didn't actually perform the experiment.'' \nb{The calculation is right, and the experiment had been performed: Hong, Ou and Mandel reported it in 1987, and the post links to it. The result does not depend on the post's choice of $i$ for deflection. The bookkeeping leaves out a factor of $\sqrt{2}$ on each two-photon amplitude; here it is the same for both detectors and changes nothing.}

The cancellation happens only because a configuration says ``a photon here, a photon there,'' not which photon is which. If configurations tracked the photons, the two routes would lead to different configurations, and one-each outcomes would happen half the time. The same holds for ``an electron here, an electron there.'' \nb{For electrons the two routes combine with opposite signs, so in this setup two electrons always leave by different sides; this was observed in 2013. Photons must also match in polarization, frequency and timing, not only in place.}

With probabilities you can merge or split possibilities as you like, because the probability that one of two exclusive events happens is the sum of their probabilities. Squared moduli do not add that way: the squared modulus of $(2+i)+(1-i)$ is 9, while the parts give 5 and 2. If the rule were linear, quantum physics would be replaced by ``a classical physics.'' \nb{A 2004 result of Scott Aaronson's points the same way.}

So ``If configurations were just states of knowledge, you could reorganize them however you liked.'' They are ``nailed in place.'' \nb{The experiment rules out photons with labels and ordinary probabilities. Whether amplitudes are physical is a question of interpretation; see ``Configurations and Amplitude.''}

Physics cannot be divided into particles with their own identities. People find quantum physics hard, in part, because they keep picturing billiard balls. ``Ha ha! Silly humans.''
